\documentclass[10pt]{article}
\usepackage[margin=20mm]{geometry}
\usepackage{amsmath,amssymb,mathtools,fontspec,unicode-math,xurl,hyperref}
\IfFontExistsTF{DejaVu Serif}{\setmainfont{DejaVu Serif}}{\setmainfont{Latin Modern Roman}}
\setmathfont{latinmodern-math.otf}
\setlength{\parindent}{0pt}\setlength{\parskip}{6pt}\sloppy
\title{Discrete trace seminorms, lifting, and the range of the geometric construction}
\author{AI-assisted technical review}
\date{}
\begin{document}
\maketitle
Abstract: The local and nonlocal components of the proposed trace seminorm have the expected common kernel and scaling. The principal Hardy, counting, and averaging mechanisms support the trace and lifting estimates under the stated regularity and quasi-uniformity hypotheses. An exact one-cell calculation also confirms the Schur-complement and boundary-kernel construction used in the numerical section. The printed cube proof nevertheless selects points outside the cube in dimensions at least three. A fixed reduction of the path height repairs this geometric step without changing the powers in the estimates. The extension to spaces without independent face values needs its trace range specified, and the local-to-global argument needs a clear boundary-connectivity convention. The assessment is Minor, with these qualifications and the limits of the review made explicit.\par
\section*{Source and precise setting}
We review S. Badia, J. Droniou, and J. Tushar, A Discrete Trace Theory for Non-Conforming Polytopal Hybrid Discretisation Methods, Foundations of Computational Mathematics 26 (2026), 2667–2716, DOI 10.1007/s10208-025-09734-6 [paper]. The acquired 50-page file is the publisher version of record, showing online publication on 3 December 2025. This review concerns that file. The one-based PDF page \(p\) corresponds to journal page \(2666+p\).\par
Section 2, PDF pp. 4–6, assumes a polytopal domain in \(\mathbb R^d\), \(d\ge2\), a regular mesh in the sense of its cited reference, and quasi-uniformity \(h\le\Psi h_t\) with \(\Psi\) independent of \(h\). These are substantive hypotheses. The results do not assert uniformity over arbitrarily graded meshes, degenerating mesh regularity, or unbounded polynomial degrees. The hidden constants may depend on the domain, mesh regularity, quasi-uniformity parameter, and polynomial degrees.\par
For the independent cell and face polynomial unknowns, write \[E_h(v)=\sum_t\left(\|\nabla v_t\|_{L^2(t)}^2+
       \sum_{f\subset\partial t}h_t^{-1}
            \|v_f-v_t\|_{L^2(f)}^2\right),\] and, with \(m_f\) the average of \(w_f\), \[T_h(w)=\sum_{f\subset\partial\Omega}
      h_f^{-1}\|w_f-m_f\|_{L^2(f)}^2
 +\sum_{f\ne g}|f|\,|g|\frac{|m_f-m_g|^2}{|x_f-x_g|^d}.\] The second sum uses ordered pairs. Theorems 2.2–2.3 assert \(T_h(\gamma v)\lesssim E_h(v)\) and a linear right inverse \(L_h\) with \(E_h(L_hw)\lesssim T_h(w)\). Both \(E_h\) and \(T_h\) are squared seminorms.\par
All 50 pages were read as extracted text, with detailed mathematical checks concentrated on the definitions, central counting and Hardy steps, lifting weights, localisation, and finite-dimensional spectral construction. This does not certify every external lemma or every geometric case. The source’s external mesh-regularity definition and auxiliary discrete functional-analysis results were not independently re-proved. The recorded bounded acquisition search found no standalone supplement. No author software or numerical experiment was run; the finite-dimensional checks below are exact algebra.\par
\section*{Kernel, scaling, and the main estimates}
Every term in \(T_h\) is nonnegative. Its vanishing forces each \(w_f\) to be constant and every pair of face averages to agree. Thus its kernel is the single global constant mode. On a connected face-adjacency mesh, \(E_h=0\) likewise makes each cell constant and transfers that constant across shared face values. The right-inverse claim is compatible with the kernels.\par
Under \(x=\mu\widehat x\) with unchanged field values, both squared forms scale by \(\mu^{d-2}\): the gradient term contributes \(\mu^{d-2}\), the penalty contributes \(\mu^{-1}\mu^{d-1}\), and the pair term contributes \(\mu^{2d-2}/\mu^d\). The seminorms themselves scale by \(\mu^{(d-2)/2}\). Remark 2.4 on PDF p. 6 should refer to squared seminorms for its exponent \(d-2\). Its conclusion about unchanged comparison constants under common rescaling is correct.\par
The nonlocal trace estimate groups face pairs in shells of radius approximately \(\ell h\). Lemma 3.2, PDF p. 14, gives at most \(C\ell^{d-2}\) faces in a shell and bounded overlap of vertical cell columns. The factor in (3.31), PDF p. 15, is \[\frac{h^{d-1}h^{d-1}h^{2-d}}{(\ell h)^d}=\ell^{-d}.\] Shell counting leaves \(\ell^{-2}\), before taking the square root. With \(r_m^2=\sum_{t\in L_m}|v_t|_{1,t}^2\), the vertical estimate becomes \[\zeta_\ell\lesssim\frac1\ell\sum_{m=0}^{\ell}r_m.\] Lemma 3.3, stated on PDF p. 15 and proved on pp. 20–21, controls the squared sum by \(18\sum_mr_m^2\). Its summation by parts and Young inequality are consistent: \[\sum_{\ell=1}^L R_\ell^2
 \le\frac4\alpha\sum_{\ell=1}^L R_\ell^2
       +(\alpha+1)\sum_{m=0}^Lr_m^2.\] Taking \(\alpha=8\) gives \(18\), including the \(m=0\) layer. The horizontal counting in Lemma 3.4, PDF p. 17, cancels the remaining shell powers; bounded layer overlap on PDF p. 18 completes that estimate, subject to the geometric correction below.\par
The lifting (4.13), PDF p. 26, averages boundary face means using \(\rho_t(f)=|f|/|A_t|\) on \(A_t\), averages neighboring cell values at interior faces, and retains the prescribed boundary polynomials. It is linear, preserves constants, and satisfies \(\gamma L_hw=w\) exactly. Its cell gradients vanish. Normalization \(\sum_f\rho_t(f)=1\) permits subtraction of a reference mean in (4.15)–(4.17), PDF p. 27. This cancellation is essential. Lemmas 4.1–4.3, PDF pp. 28–29, give \[\sum_fD_g\rho(f)\lesssim h/\delta_g.\] For faces separated by \(r\), the interior-overlap contribution in Lemma 4.2, PDF pp. 33–34, has the correct tail \(h\sum_{\delta\gtrsim r}\delta^{-d-1}\lesssim r^{-d}\). The symmetric-difference contribution has bounded multiplicity and the same pair weight. Near-grid-scale terms must be retained separately when using the integral comparison; its lower endpoint cannot be moved past contributing layers.\par
As a further exact check, if every face mean is zero, all cell and interior-face values produced by the lift vanish. Its energy then consists only of \(\sum_{f\subset\partial\Omega}h_t^{-1}\|w_f\|_{L^2(f)}^2\), controlled by the local trace term under mesh regularity. No boundary polynomial degree is lost.\par
For the global lifting, PDF pp. 40–44, averaged partition functions preserve the trace because \(\sum_i\pi_f^0\eta_i=1\). The cutoff differences are \(O(h_t)\), generating the bulk \(L^2\) term in (5.19). The local bulk estimate and global zero-mean boundary Poincaré inequality provide the stated control. This argument uses the full boundary seminorm, so the separate local-seminorm connectivity concern below does not itself refute the global lifting.\par
\section*{Geometric and range corrections}
Problem 1 (Minor: the height used in the cube proof).  Section 3.2, PDF pp. 9–10, considers a unit cube and places the horizontal path for a pair \(f,g\) at height \(r=|x_f-x_g|\). Section 3.3.2, PDF p. 12, requires cells containing its raised endpoints for every pair of side faces. For \(d\ge3\), these endpoints need not belong to the cube.\par
Take \(\Omega=(0,1)^3\) with a uniform \(4\times4\times4\) Cartesian mesh. On the bottom side choose \[x_f=(1/8,1/8,0),\qquad x_g=(7/8,7/8,0).\] This is regular and quasi-uniform, but \(r=3\sqrt2/4>1\). Both raised endpoints lie above the cube. The selected cells do not exist and the path-cell set (3.17) is empty. The decomposition on PDF p. 12 cannot be made as printed. The graph-following description on PDF p. 37 again uses the full horizontal distance as height.\par
A direct repair uses height \(\kappa r\) with \(\kappa=1/(2\sqrt{d-1})\), keeping every raised path below height \(1/2\). Vertical telescoping reaches layers \(m\le\kappa\ell+O(1)\); the Hardy estimate applies, with the bounded endpoint terms estimated separately. Horizontal shell and cone counts are unchanged because their projections are unchanged. Path cells now lie in layers \(|m-\kappa\ell|\le O(1)\). For fixed \(m\), only \(O(1+\kappa^{-1})\) shell indices contribute. The overlap remains bounded independently of \(h\), and the powers in the estimates are preserved. A revised proof should state the adjusted layer ranges and choose nested boundary charts with adequate vertical clearance.\par
This is an undefined intermediate construction, not a counterexample to the trace inequality. The repair changes dimension-dependent constants rather than the asserted order, supporting a Minor proof correction.\par
Problem 2 (Minor qualification: spaces without face values).  Remark 5.7, PDF pp. 44–45, embeds broken cell-polynomial spaces by assigning averaged interior traces and actual boundary traces. An embedding transfers a trace bound, but does not supply a right inverse onto the full product of independent boundary-face polynomials.\par
For a single square cell with degree-zero DG space, a scalar cell value \(c\) induces \((c,c,c,c)\). The independent face datum \((1,0,0,0)\) has no such lift. The hybrid construction’s independent boundary values cannot simply be interpreted as a DG function. This exact range obstruction occurs on a regular mesh.\par
A precise adaptation must specify the actual DG trace image as target and justify a lifting into the cell space with the appropriate energy estimate, or enrich the space if independent face data are intended. If the remark intends only that natural image, the example is outside its target. The issue is therefore a qualification of the adaptation, not a failure of the hybrid theorem. The mixed-degree argument on PDF p. 44 does work: the constructed cell and interior-face values are constants and the prescribed boundary polynomials retain their allowed degree.\par
A further domain convention should be explicit. Proposition 5.2, PDF pp. 38–40, joins any two boundary charts by a chain of nonempty boundary overlaps. This uses connectedness of the boundary-overlap graph. If polytopal domain includes a polygonal annulus, take \(\Omega=(-2,2)^2\setminus[-1,1]^2\), with charts each meeting only one boundary component. Data zero on the outer component and one on the inner component have every local seminorm zero, but \[T_h(w)=2\sum_{\substack{f\ \mathrm{outer}\\g\ \mathrm{inner}}}
              \frac{|f|\,|g|}{|x_f-x_g|^2}>0.\] The local-seminorm equivalence needs component-mean terms in this setting. If the intended convention for polytope excludes holes, the example is outside the intended class. The terminology alternates in the text, so this is a hypothesis clarification, not an adopted counterexample to the main theorems.\par
\section*{An exact spectral check and the assessment}
Section 6, PDF pp. 45–47, compares the boundary seminorm with the Schur complement of the bulk form after a rank-one kernel correction. A one-cell calculation checks its normalization independently of any experiment.\par
Let the cell be the unit square, \(k=0\), and label face values \(w_1,\ldots,w_4\) cyclically. As \(h_t=\sqrt2\) and each face has length one, \[E(v,w)=\frac1{\sqrt2}\sum_{i=1}^4(w_i-v)^2,\qquad
 v_{\min}=\frac14\sum_iw_i.\] Writing \(J=ee^T\), where \(e=(1,1,1,1)^T\), the Schur matrix is \[A_{\rm SC}=\frac1{\sqrt2}(I-J/4).\] Adjacent face centers have squared distance \(1/2\), opposite centers squared distance \(1\). Let \(L\) be the graph Laplacian with adjacent weights \(2\) and opposite weights \(1\). The ordered-pair trace matrix is \(H=2L\), with eigenvalues \(0,12,12,16\). For example, \(w=(1,-1,1,-1)\) has trace energy \(64\) and minimized bulk energy \(2\sqrt2\). The vector \((1,0,-1,0)\) gives \(24\) and \(\sqrt2\), respectively.\par
Here \(S=(1,1,1,1)\), so \(S^TS=J\). The generalized eigenvalues of \((A_{\rm SC}+J,H+J)\) are \[1,\qquad \frac1{12\sqrt2},\qquad
 \frac1{12\sqrt2},\qquad \frac1{16\sqrt2}.\] The constant mode gives \(1\), while the other modes have zero boundary mean. This supports the pair ordering, Schur construction, and kernel correction. It is a fixed-mesh check, not proof of an \(h\)-uniform bound.\par
Problem 3 (Minor: the printed local assembly formula).  On PDF p. 46 the local array is defined through differences \(\Delta_{ij}=\overline\phi_i^{\,f_l}-\overline\phi_j^{\,f_m}\) of scalar basis-function means and a contribution \(K\Delta^T\Delta\). This does not give the face-pair energy in (2.2) under its literal displayed interpretation.\par
For \(k=0\), choose the constant basis \(1\) on each unit-length face. Every basis mean equals one, hence every entry of \(\Delta\) is zero. The printed nonlocal contribution is zero, as is the local oscillation contribution. In contrast, the pair with coefficients \((c_l,c_m)=(1,0)\) has energy \(K(c_l-c_m)^2=K>0\).\par
Let \(a_l,a_m\) be the vectors of basis means on the two faces. The correct pair vector is the concatenation \(b=(a_l^T,-a_m^T)^T\), with pair matrix \(Kbb^T\). Its quadratic form is exactly \(K(a_l^Tc_l-a_m^Tc_m)^2\). For constants this gives \(K\left(\begin{smallmatrix}1&-1\\-1&1\end{smallmatrix}\right)\). The assembly must also retain the ordered-pair convention. An explicit embedding into the combined pair-coordinate space could clarify the intended notation. The theory does not depend on this paragraph, and the reported nondegenerate \(k=0\) spectra are counterevidence against assuming the program literally used the zero formula. No actual-code defect or incorrect reported spectrum is inferred.\par
Figure 8’s finite Cartesian-square spectra provide illustrations; they do not verify all dimensions and regular polytopal meshes. No discrepancy in the plotted values is claimed. The review supports the proposed seminorm and the inspected analytic mechanisms, with the concrete path-height repair and the specified trace-range and domain qualifications, and a correction to the printed pair-matrix formula. It does not establish degree-robust constants, a graded-mesh extension, or an independently reproduced numerical result. No supported Major defect in the main estimates was established.\par
1 S. Badia, J. Droniou, and J. Tushar, A Discrete Trace Theory for Non-Conforming Polytopal Hybrid Discretisation Methods, Found. Comput. Math. 26 (2026), 2667–2716. Published online 3 December 2025. https://doi.org/10.1007/s10208-025-09734-6.\par
\end{document}
